%% kapa: exam=1 round=37 year=2026 period=1 form=A subject=economics subject_ko=경제학원론 seq=5 number=45
%% source: Q-Net 감정평가사 기출문제 (cst00309) artlSeq=5258924 file=raw/1008/5258924/2026년 제37회 감정평가사 1차 1교시.pdf
%% conversion: kordoc 4.24.51 -> tools/gen1.py (auto); review=proofread-sonnet; flags=none
\begin{question}{45}
콥-더글라스 생산함수 $Q=L^{\frac{1}{3}}K^{\frac{2}{3}}$에 관한 설명으로 옳지 않은 것은? (단, 재화 및 생산요소의 가격은 일정하며, $Q$는 생산량, $L$은 노동, $K$는 자본이다.)
\begin{choices}
  \item 규모에 대한 수익불변이다.
  \item 노동에 대해 한계생산체감의 법칙이 작용한다.
  \item 완전경쟁 노동시장에서 노동의 소득분배율(노동소득/총생산)은 1/3이다.
  \item 1차 동차 함수이다.
  \item 한계기술대체율($MRTS_{LK}$)은 $2\dfrac{K}{L}$이다.
\end{choices}
\end{question}
\answer{5}
